\section*{5 Proof of Theorem1}
As before, let $\mathtt C$ be a depth four circuit computing the polynomial $f_n$. Before proving the theorem, let us first justify the assumption that $\mathbb F$ is an algebraically closed field that we made in Section3. Suppose not. Then, let $\overline{\mathbb F}$ be its algebraic closure. Since $\mathtt C$ is also a circuit over $\overline{\mathbb F}$ and $f_n$ a polynomial over $\overline{\mathbb F}$, we can make all arguments assuming the underlying field to be $\overline{\mathbb F}$. Since the size of a circuit does not depend on the underlying field, the lower bound so obtained will continue to hold when we treat $\mathtt C$ as a circuit over $\mathbb F$.
First we will prove a lower bound on the number of wires of $\mathtt C$. If the number of monomials in $\mathtt C$ is $\lfloor m^2d_x/(\ln m)^5\rfloor$ then there is nothing to prove. Otherwise from Lemma10, we can obtain a circuit $\mathtt C_1$ such that the support of all the monomials of $\mathtt C_1$ is at most $\tau=\lfloor20\ln m\rfloor$, the number of gates in $\mathtt C_1$ is at most the number of gates in $\mathtt C$ and the fan-in of each gate in $\mathtt C_1$ is upper bounded by the fan-in of the corresponding gate in $\mathtt C$. Then, if $\mathtt C_1$ does not satisfy the hypothesis of Lemma12, the size of $\mathtt C_1$ and hence the size of $\mathtt C$ is at least $\Omega(m^2d_x/(\ln m)^5)$. Otherwise, we can obtain a pruned circuit $\mathtt D$ such that the top fan-in and the bottom support of $\mathtt D$ are upper bounded by the top fan-in and bottom support of $\mathtt C_1$ and so proving a lower bound on the top fan-in of $\mathtt D$ would suffice.
As $\mathtt D$ computes $f_1$ (defined in Section4), $\PSP_{M,k,\ell}(\mathtt D)=\PSP_{M,k,\ell}(f_1)$. Lemma7 and 16 imply
\begin{equation}
s\ge\frac{m^{-O(1)}\min\{4^{-k}\binom m\ell\binom mk,\binom m{\ell+d_x-k}\}}
{m^{O(1)}\binom m{\ell+2kt\tau}\binom{\lceil w/t\rceil+k-1}k},\tag{1}
\end{equation}
and the required lower bound follows from the next claim, which is proved in AppendixC.
\begin{authorclaim}{Claim17}
The top fan-in $s$ of $\mathtt D$ is $\omega(m^2d_x/(\ln m)^5)$.
\end{authorclaim}
Now let us prove the lower bound on the number of gates. Notice that if the circuit $\mathtt C$ computing $f$ has a heavy gate as defined in Section3 then we are done. So assume that it does not have any heavy gates. Now, if the number of monomials in $\mathtt C$ is $\lfloor m^2d_x/(\ln m)^5\rfloor$ then there is nothing to prove. Otherwise from Lemma10, we can obtain a circuit $\mathtt C_1$ such that the support of all the monomials of $\mathtt C_1$ is at most $\tau=\lfloor20\ln m\rfloor$, the number of gates in $\mathtt C_1$ is at most the number of gates in $\mathtt C$ and the fan-in of each gate in $\mathtt C_1$ is upper bounded by the fan-in of the corresponding gate in $\mathtt C$. Obtain a circuit $\mathtt D$ from $\mathtt C_1$ by picking a set $M_2\subseteq[3m]\setminus M_1$ (where $M_1$ is as in Lemma10), $|M_2|=m$ arbitrarily and setting variables in $\mathbf x_{M_2}$ and $\mathbf y_{M_2}$ to 0 (notice that the top fan-in and bottom support of $\mathtt D$ are upper bounded by the top fan-in and bottom support of $\mathtt C_1$). Now $\mathtt D$ computes $f_1$ (where $f_1$ is as defined in Section4) and just as it was done in the preceding paragraph, we can show that the top fan-in of $\mathtt D$ is $\omega(m^2d_x/(\ln m)^5)$. However, we only get an $\Omega(md_x/(\ln m)^3)$ lower bound on the number of gates since the definition of a heavy gate is the bottleneck.
